\section{Numerical resolution and convergence}
\label{app:repro}

\paragraph{Resolution.}  The reported solutions use polynomial degree $n=22$, that is
$N=23$ Legendre modes per field ($N-2=21$ for the grounded potential), and every integral is evaluated
with $\max(64,2n+10)$ Gauss--Legendre nodes, which resolves the mass, stiffness
and coupling matrices to machine precision at the degrees used here.  The two
electrical cases differ only in the potential space: the grounded basis is built
from the combinations that vanish on both faces, and the open-circuit basis
keeps all $N$ modes, because a constant potential is not a gauge for a nonzero in-plane wavenumber.  No field is discretized in $x$ or in $y$.

\paragraph{Convergence and conservation.}  Between degree $18$ and degree $30$
both the operating-point share and the branch maximum change by less than
$10^{-5}$, and successive differences at the branch maximum decrease
monotonically with degree.  Refining the quadrature from $64$ to $256$ nodes at
fixed basis changes no reported digit.  The energy identity
\eqref{eq:energy_identity} is satisfied to $4\times10^{-12}$ of the incident
power at the reported solutions, and the partition at the branch maximum varies by $1.0\times10^{-6}$ in absolute terms between
degrees $14$ and $34$.  These are the quantities that
license quoting the shares to the precision used in the text.
